\documentclass[aspectratio=169,11pt]{beamer}

\usepackage{beamerthemeSWUFE}
\usepackage{graphicx}
\usepackage{booktabs}
\usepackage{amsmath}
\usepackage{amssymb}
\usepackage{listings}
\usepackage{pgfplots}
\pgfplotsset{compat=1.18}
\graphicspath{{assets/figures/}{assets/logos/}}

\lstdefinestyle{swufeCode}{
  basicstyle=\ttfamily\scriptsize,
  keywordstyle=\color{swufeBlue}\bfseries,
  commentstyle=\color{swufeGray},
  stringstyle=\color{swufeGold},
  numbers=left,
  numberstyle=\tiny\color{swufeGray},
  frame=single,
  rulecolor=\color{swufeLine},
  breaklines=true,
  columns=fullflexible,
  xleftmargin=1.5em,
  framexleftmargin=1.1em
}
\lstset{style=swufeCode}

\title{本科毕业论文题目示例}
\subtitle{西南财经大学本科毕业论文答辩模板\\SWUFE Undergraduate Thesis Defense Beamer Template}
\author{学生姓名}
\studentid{202200000000}
\school{学院名称}
\major{专业名称}
\supervisor{指导教师}
\date{20XX年X月}

\begin{document}

\begin{frame}[plain,noframenumbering]
  \titlepage
\end{frame}

\begin{frame}{目录}
  \tableofcontents
\end{frame}

\section{研究背景与意义}

\begin{frame}{研究背景}{说明选题来源与现实问题}
  \begin{itemize}
    \item 简要说明研究主题所处的行业、市场、政策或学术背景。
    \item 指出现实中仍然存在的主要问题，例如效率不足、风险识别不充分、评价体系不完善或管理机制有待优化。
    \item 说明本文为什么值得研究：可以来自理论补充、实务应用、数据证据或政策参考价值。
  \end{itemize}
  \begin{alertblock}{写作提示}
    本页用于展示普通 bullet page。请将占位文字替换为自己的研究背景，避免堆砌过长段落。
  \end{alertblock}
\end{frame}

\begin{frame}{研究意义}{三点式表达}
  \begin{swufekeypoints}
    \swufepoint{理论意义}{01}{梳理已有研究脉络，明确本文在概念、机制或解释框架上的补充空间。}
    \swufepoint{方法意义}{02}{构建适合本科论文的研究设计，包括变量定义、样本选择、模型设定或案例分析框架。}
    \swufepoint{实践意义}{03}{结合研究结果提出可理解、可落地的启示，为相关主体的决策提供参考。}
  \end{swufekeypoints}
\end{frame}

\section{文献综述与理论基础}

\begin{frame}{文献综述}{从主题到研究缺口}
  \begin{columns}[T,onlytextwidth]
    \begin{column}{0.55\textwidth}
      \begin{itemize}
        \item 第一类文献：围绕核心概念、制度背景或理论机制展开。
        \item 第二类文献：使用经验数据、案例材料或比较研究方法进行检验。
        \item 第三类文献：讨论研究对象的影响因素、经济后果或治理启示。
        \item 研究缺口：现有研究在样本、方法、视角或解释深度方面仍有拓展空间。
      \end{itemize}
    \end{column}
    \begin{column}{0.4\textwidth}
      \begin{block}{综述逻辑}
        核心概念界定\\
        $\Downarrow$\\
        理论机制梳理\\
        $\Downarrow$\\
        经验研究比较\\
        $\Downarrow$\\
        本文研究切入点
      \end{block}
    \end{column}
  \end{columns}
\end{frame}

\begin{frame}{理论基础}{分析框架示例}
  本科论文可使用简洁的理论框架说明变量之间的关系。设研究对象的结果变量为 $Y$，核心解释变量为 $X$，控制变量为 $Z$，可写为：
  \[
    Y_i = \alpha + \beta X_i + \gamma Z_i + \varepsilon_i.
  \]
  若采用机制分析，可进一步说明：
  \[
    X \rightarrow M \rightarrow Y,
  \]
  其中 $M$ 表示可能的中介机制或传导路径。
  \begin{block}{模板提示}
    公式页不需要复杂化。答辩中应重点解释变量含义、假设方向和模型服务于哪个研究问题。
  \end{block}
\end{frame}

\section{研究方法与设计}

\begin{frame}{研究问题与假设}{将论文问题转化为可检验命题}
  \begin{enumerate}
    \item 研究问题一：研究对象的基本特征是什么？
    \item 研究问题二：核心因素是否会影响研究结果？
    \item 研究问题三：影响是否在不同样本、情境或条件下存在差异？
  \end{enumerate}
  \begin{alertblock}{假设写法示例}
    H1：在控制其他因素后，核心解释变量 $X$ 与结果变量 $Y$ 存在显著关系。该表述为模板占位，请根据真实论文调整。
  \end{alertblock}
\end{frame}

\begin{frame}[fragile]{研究流程}{代码或流程片段展示}
  \begin{lstlisting}[language=Python]
def thesis_workflow(raw_data):
    """Illustrative pseudo-code for a thesis defense template."""
    data = clean_data(raw_data)
    variables = construct_variables(data)
    baseline = estimate_baseline_model(variables)
    robustness = run_robustness_checks(variables)
    return summarize_results(baseline, robustness)
  \end{lstlisting}
  \begin{itemize}
    \item 本模板使用 `listings` 展示代码，不需要 `-shell-escape`。
    \item 若论文不涉及编程，可将本页替换为流程图、研究框架图或案例分析步骤。
  \end{itemize}
\end{frame}

\section{数据来源与样本说明}

\begin{frame}{数据来源}{样本与变量说明}
  \begin{table}
    \centering
    \caption{数据与变量说明示例（illustrative）}
    \begin{tabular}{lll}
      \toprule
      类别 & 内容 & 说明 \\
      \midrule
      样本范围 & 20XX--20XX 年 & 根据论文主题确定 \\
      研究对象 & 企业、城市、基金或问卷样本 & 替换为实际对象 \\
      结果变量 & $Y$ & 衡量研究结果或绩效表现 \\
      核心变量 & $X$ & 衡量主要解释因素 \\
      控制变量 & $Z$ & 控制规模、地区、年份等因素 \\
      \bottomrule
    \end{tabular}
  \end{table}
  \vspace{-0.3em}
  \scriptsize 注：本表为模板示例，正式答辩请替换为真实数据来源和变量定义。
\end{frame}

\begin{frame}{实证或案例设计}{研究设计的基本要素}
  \begin{columns}[T,onlytextwidth]
    \begin{column}{0.48\textwidth}
      \begin{block}{量化研究}
        \begin{itemize}
          \item 样本筛选规则
          \item 变量定义与度量
          \item 基准模型与控制变量
          \item 稳健性与异质性检验
        \end{itemize}
      \end{block}
    \end{column}
    \begin{column}{0.48\textwidth}
      \begin{block}{案例研究}
        \begin{itemize}
          \item 案例选择依据
          \item 分析维度与资料来源
          \item 关键事实与时间线
          \item 经验总结与启示
        \end{itemize}
      \end{block}
    \end{column}
  \end{columns}
\end{frame}

\section{研究结果与分析}

\begin{frame}{主要结果}{核心结果表}
  \begin{table}
    \centering
    \caption{基准结果展示示例值（illustrative）}
    \begin{tabular}{lrrrr}
      \toprule
      变量 & 模型一 & 模型二 & 模型三 & 模型四 \\
      \midrule
      核心变量 $X$ & 0.120 & 0.115 & 0.098 & 0.102 \\
      控制变量 & 否 & 是 & 是 & 是 \\
      年份固定效应 & 否 & 否 & 是 & 是 \\
      个体固定效应 & 否 & 否 & 否 & 是 \\
      样本量 & 1,000 & 1,000 & 1,000 & 1,000 \\
      \bottomrule
    \end{tabular}
  \end{table}
  \begin{alertblock}{说明}
    所有数字均为示例值 / illustrative，仅用于展示表格版式，不代表任何真实研究结论。
  \end{alertblock}
\end{frame}

\begin{frame}{双栏图文页}{图表与解释并列展示}
  \begin{columns}[T,onlytextwidth]
    \begin{column}{0.56\textwidth}
      \centering
      \swufeplaceholderfigure{研究结果趋势图占位 / illustrative figure}
    \end{column}
    \begin{column}{0.4\textwidth}
      \begin{itemize}
        \item 图表应直接服务于核心结论，避免展示与论文问题无关的信息。
        \item 解释时先说现象，再说明原因或机制，最后回到研究假设。
        \item 正式论文中请替换为由实证分析、调研或案例材料生成的真实图表。
      \end{itemize}
    \end{column}
  \end{columns}
\end{frame}

\begin{frame}{稳健性检验}{结果可靠性说明}
  \begin{table}
    \centering
    \caption{稳健性检验示例值（illustrative）}
    \begin{tabular}{lrrr}
      \toprule
      检验方式 & 核心系数 & 方向是否一致 & 结论 \\
      \midrule
      替换变量度量 & 0.105 & 是 & 支持 \\
      调整样本范围 & 0.097 & 是 & 支持 \\
      增加控制变量 & 0.101 & 是 & 支持 \\
      分组样本检验 & 0.088 & 基本一致 & 部分支持 \\
      \bottomrule
    \end{tabular}
  \end{table}
  \scriptsize 注：本页展示稳健性检验页的常见结构，数值均为示例。
\end{frame}

\begin{frame}{结果讨论}{避免过度解释}
  \begin{itemize}
    \item 如果结果支持研究假设，应说明证据来自哪些模型、表格或案例事实。
    \item 如果部分结果不显著，应如实说明，并讨论样本、变量或方法上的可能原因。
    \item 不建议使用“完全证明”“必然导致”等绝对化表述，本科论文更适合使用审慎、证据导向的语言。
  \end{itemize}
\end{frame}

\section{主要结论与不足}

\begin{frame}{主要结论}{回扣研究问题}
  \begin{enumerate}
    \item 结论一：概括研究对象的主要特征或现实问题。
    \item 结论二：说明核心变量、机制或案例事实与研究结果之间的关系。
    \item 结论三：总结稳健性、分组分析或案例比较带来的补充发现。
  \end{enumerate}
  \begin{block}{答辩建议}
    结论页应回扣研究问题、方法和证据，不宜新增复杂结果或临时扩展观点。
  \end{block}
\end{frame}

\begin{frame}{研究不足与展望}
  \begin{itemize}
    \item 数据方面：样本区间、样本规模或数据可得性可能限制结论外推。
    \item 方法方面：变量度量、模型设定或案例材料仍存在进一步改进空间。
    \item 研究方面：未来可扩展到更长时间、更细分样本或更多比较对象。
  \end{itemize}
  \begin{alertblock}{展望}
    后续研究可结合新的数据来源、理论视角或实践场景，对本文结论进行进一步验证。
  \end{alertblock}
\end{frame}

\section{致谢与答辩 Q\&A}

\begin{frame}{致谢}
  \begin{center}
    \vspace{0.7cm}
    {\LARGE\bfseries\color{swufeBlue}谢谢各位老师的指导与聆听}\\[0.8cm]
    {\Large Questions \& Answers}\\[0.5cm]
    \rule{0.55\linewidth}{0.6pt}\\[0.5cm]
    \swufeUniversityCN\\
    \swufeUniversityEN
  \end{center}
\end{frame}

\end{document}
